\chapter{Droites, cercles et angles}\label{ch:g6:lines}

La géométrie commence avec une règle, une équerre et un compas. Ce
chapitre fixe le vocabulaire --- points, \omterm{def:g6:lines:objects}{segments}, demi-droites,
droites, cercles ---, introduit les deux positions spéciales des
droites (\omterm{def:g6:lines:perp}{parallèles} et \omterm{def:g6:lines:perp}{perpendiculaires}), et enseigne à mesurer et
tracer des angles.

\section{Points, segments, droites}

\begin{definition}[Objets de base]\label{def:g6:lines:objects}
Par deux points distincts $A$ et $B$ passent~:
\begin{itemize}
  \item le \emph{segment}\index{segment} $[AB]$~: la partie de la
        droite entre $A$ et $B$ (il a une longueur, notée $AB$)~;
  \item la \emph{demi-droite} $[AB)$~: part de $A$, passe par $B$ et
        continue indéfiniment~;
  \item la \emph{droite}\index{droite} $(AB)$~: s'étend indéfiniment
        des deux côtés. Deux points déterminent exactement une
        droite.
\end{itemize}
Des points sur la même droite sont dits \emph{alignés}.
\end{definition}

\begin{omfigure}
\begin{tikzpicture}[scale=1.0]
  \fill (0,1.8) circle (1.5pt) node[above, font=\small] {$A$};
  \fill (2.5,1.8) circle (1.5pt) node[above, font=\small] {$B$};
  \draw[thick, omDef] (0,1.8) -- (2.5,1.8);
  \node[right, font=\small] at (2.9,1.8) {segment $[AB]$};
  \fill (0,0.9) circle (1.5pt) node[above, font=\small] {$A$};
  \fill (2.5,0.9) circle (1.5pt) node[above, font=\small] {$B$};
  \draw[thick, omThm] (0,0.9) -- (4.1,0.9);
  \node[right, font=\small] at (4.3,0.9) {demi-droite $[AB)$};
  \fill (0,0) circle (1.5pt) node[above, font=\small] {$A$};
  \fill (2.5,0) circle (1.5pt) node[above, font=\small] {$B$};
  \draw[thick, omMeth] (-1.4,0) -- (4.1,0);
  \node[right, font=\small] at (4.3,0) {droite $(AB)$};
\end{tikzpicture}

{\small Mêmes deux points, trois objets différents. Les crochets
disent ce qui s'arrête et ce qui continue~: $[$ s'arrête, $($
continue.}
\end{omfigure}

\section{Droites parallèles et perpendiculaires}

\begin{definition}[Perpendiculaires, parallèles]\label{def:g6:lines:perp}
Deux droites sont \emph{perpendiculaires}\index{droites perpendiculaires}\index{perpendiculaire}
lorsqu'elles se croisent à \omterm{def:g3:shapes:rightangle}{angle droit}~; on écrit $d_1 \perp d_2$ et on
marque l'\omterm{def:g3:shapes:rightangle}{angle droit} d'un petit carré. Deux droites sont
\emph{parallèles}\index{droites parallèles}\index{parallèle} lorsqu'elles ne se
rencontrent jamais, aussi loin qu'on les prolonge~; on écrit
$d_1 \parallel d_2$.
\end{definition}

\begin{omfigure}
\begin{tikzpicture}[scale=1.0]
  \draw[thick] (-1.8,0) -- (1.8,0) node[right, font=\small] {$d_1$};
  \draw[thick, omDef] (0.6,-1.2) -- (0.6,1.4)
    node[above, font=\small] {$d_2$};
  \draw[thin] (0.6,0) ++(-0.25,0) -- ++(0,0.25) -- ++(0.25,0);
  \node[below, font=\small] at (0,-1.5) {$d_1 \perp d_2$};
  \begin{scope}[xshift=6.5cm]
  \draw[thick] (-1.8,-0.5) -- (1.8,0.3) node[right, font=\small] {$d_3$};
  \draw[thick, omThm] (-1.8,-1.3) -- (1.8,-0.5)
    node[right, font=\small] {$d_4$};
  \node[below, font=\small] at (0,-1.5) {$d_3 \parallel d_4$};
  \end{scope}
\end{tikzpicture}

{\small Un croisement à \omterm{def:g3:shapes:rightangle}{angle droit} (marqué du petit carré) et deux
\omterm{def:g6:lines:perp}{parallèles}~: même direction, pas de point de croisement.}
\end{omfigure}

\begin{proposition}[Deux faits utiles]\label{prop:g6:lines:facts}
\begin{enumerate}
  \item Si deux droites sont toutes deux \omterm{def:g6:lines:perp}{perpendiculaires} à une
        troisième, elles sont \omterm{def:g6:lines:perp}{parallèles} entre elles.
  \item Si deux droites sont \omterm{def:g6:lines:perp}{parallèles}, toute droite \omterm{def:g6:lines:perp}{perpendiculaire}
        à l'une est \omterm{def:g6:lines:perp}{perpendiculaire} à l'autre.
\end{enumerate}
\end{proposition}
\admitted

\begin{method}[Tracer avec l'équerre]\label{met:g6:lines:setsquare}
Pour tracer la \omterm{def:g6:lines:perp}{perpendiculaire} à une droite $d$ passant par un point
$P$~:
\begin{enumerate}
  \item placer un côté de l'\omterm{def:g3:shapes:rightangle}{angle droit} de l'équerre le long de $d$~;
  \item glisser l'équerre le long de $d$ jusqu'à ce que l'autre côté
        atteigne $P$~;
  \item tracer la droite le long de ce côté, et marquer l'\omterm{def:g3:shapes:rightangle}{angle droit}.
\end{enumerate}
Pour une \omterm{def:g6:lines:perp}{parallèle} par $P$~: tracer une \omterm{def:g6:lines:perp}{perpendiculaire} à $d$, puis la
\omterm{def:g6:lines:perp}{perpendiculaire} à \emph{cette} droite par $P$
(\cref{prop:g6:lines:facts}).
\end{method}

\section{Cercles}

\begin{definition}[Cercle]\label{def:g6:lines:circle}
Le \emph{cercle}\index{cercle} de centre $O$ et de \omterm{def:g3:shapes:shapes}{rayon} $r$ est
l'ensemble de tous les points situés à distance exactement $r$ de
$O$. Un \omterm{def:g6:lines:objects}{segment} du centre au cercle est un \emph{\omterm{def:g3:shapes:shapes}{rayon}}~; un \omterm{def:g6:lines:objects}{segment}
joignant deux points du cercle en passant par le centre est un
\emph{\omterm{def:g4:geometry:circle}{diamètre}} --- sa longueur est $2r$~; un \omterm{def:g6:lines:objects}{segment} joignant deux
points du cercle est une \emph{corde}.
\end{definition}

\begin{omfigure}
\begin{tikzpicture}[scale=1.15]
  \draw[thick] (0,0) circle (1.5);
  \fill (0,0) circle (1.5pt) node[below left, font=\small] {$O$};
  \fill (1.5,0) circle (1.5pt) node[right, font=\small] {$M$};
  \draw[omDef, thick] (0,0) -- (1.5,0)
    node[midway, above, font=\small] {rayon};
  \fill (-1.06,1.06) circle (1.5pt) node[above left, font=\small] {$A$};
  \fill (1.06,-1.06) circle (1.5pt) node[below right, font=\small] {$B$};
  \draw[omThm, thick] (-1.06,1.06) -- (1.06,-1.06)
    node[pos=0.25, right=2pt, font=\small] {diamètre};
  \fill (0.52,1.41) circle (1.5pt) node[above, font=\small] {$C$};
  \fill (1.41,0.52) circle (1.5pt) node[right, font=\small] {$D$};
  \draw[omMeth, thick] (0.52,1.41) -- (1.41,0.52)
    node[midway, above right=-2pt, font=\small] {corde};
\end{tikzpicture}

{\small Un cercle avec un \omterm{def:g3:shapes:shapes}{rayon} $[OM]$, un \omterm{def:g4:geometry:circle}{diamètre} $[AB]$ (une corde
passant par le centre, deux fois plus longue que le \omterm{def:g3:shapes:shapes}{rayon}) et une
corde $[CD]$.}
\end{omfigure}

\begin{example}\label{ex:g6:lines:circle}
«~Tracer le cercle de centre $O$ passant par $A$~»~: ouvrir le
compas de $O$ à $A$ --- le \omterm{def:g3:shapes:shapes}{rayon} est la distance $OA$ --- et tourner.
Tout point de ce cercle est à la même distance de $O$ que $A$.
\end{example}

\section{Angles}

\begin{definition}[Angle]\label{def:g6:lines:angle}
Deux demi-droites $[AB)$ et $[AC)$ de même origine $A$ forment
l'\emph{angle}\index{angle} $\widehat{BAC}$~; le point $A$ est son
\emph{\omterm{def:g5:solids:def}{sommet}} (toujours la lettre du milieu~!). Les angles se mesurent
en \emph{degrés} ($^\circ$), de $0^\circ$ à $360^\circ$ pour un tour
complet. Un angle est~:
\begin{itemize}
  \item \emph{droit} s'il mesure $90^\circ$~;
  \item \emph{aigu} s'il mesure moins de $90^\circ$~;
  \item \emph{obtus} s'il mesure entre $90^\circ$ et $180^\circ$~;
  \item \emph{plat} s'il mesure $180^\circ$ (les deux demi-droites
        forment une droite).
\end{itemize}
\end{definition}

\begin{omfigure}
\begin{tikzpicture}[scale=0.9]
  \draw[thick] (0,0) -- (1.8,0);
  \draw[thick] (0,0) -- (40:1.8);
  \draw[omDef, ->] (0.7,0) arc (0:40:0.7);
  \node[omDef, font=\small] at (20:1.0) {$40^\circ$};
  \node[below, font=\small] at (0.9,-0.3) {aigu};
  \begin{scope}[xshift=3.6cm]
  \draw[thick] (0,0) -- (1.8,0);
  \draw[thick] (0,0) -- (0,1.8);
  \draw[thin] (0.3,0) -- (0.3,0.3) -- (0,0.3);
  \node[below, font=\small] at (0.9,-0.3) {droit};
  \end{scope}
  \begin{scope}[xshift=7.2cm]
  \draw[thick] (0,0) -- (1.8,0);
  \draw[thick] (0,0) -- (135:1.8);
  \draw[omThm, ->] (0.6,0) arc (0:135:0.6);
  \node[omThm, font=\small] at (65:1.0) {$135^\circ$};
  \node[below, font=\small] at (0.9,-0.3) {obtus};
  \end{scope}
  \begin{scope}[xshift=11.6cm]
  \draw[thick] (-1.5,0) -- (1.5,0);
  \fill (0,0) circle (1.2pt);
  \draw[omMeth, ->] (0.5,0) arc (0:180:0.5);
  \node[below, font=\small] at (0,-0.3) {plat};
  \end{scope}
\end{tikzpicture}

{\small Les quatre familles d'angles. L'\omterm{def:g3:shapes:rightangle}{angle droit} ($90^\circ$) est
la référence~: aigu signifie plus petit, obtus signifie plus grand.}
\end{omfigure}

\begin{method}[Mesurer un angle avec un rapporteur]\label{met:g6:lines:protractor}
\begin{enumerate}
  \item Placer le centre du rapporteur exactement sur le \omterm{def:g5:solids:def}{sommet} de
        l'angle~;
  \item aligner sa ligne $0^\circ$ avec un côté de l'angle~;
  \item lire la graduation croisée par l'autre côté --- en utilisant
        l'échelle qui commence à $0$ du côté aligné~;
  \item vérifier à l'œil~: un angle aigu doit lire moins de
        $90^\circ$, un obtus plus.
\end{enumerate}
\end{method}

\begin{example}\label{ex:g6:lines:estimate}
Avant de mesurer, estimer~! Un angle un peu plus ouvert que le coin
d'une feuille de papier est un peu plus de $90^\circ$~; la \omterm{def:g3:division:half}{moitié} d'un
\omterm{def:g3:shapes:rightangle}{angle droit} est $45^\circ$~; le tiers d'un \omterm{def:g3:shapes:rightangle}{angle droit} est $30^\circ$.
Si le rapporteur dit $150^\circ$ pour un angle qui semble aigu, on a
lu la mauvaise échelle~: la mesure correcte est
$180^\circ - 150^\circ = 30^\circ$.
\end{example}

\section{Exercices}

\begin{exercise}[$\star$]\label{exo:g6:lines:1}
Tracer trois points $A$, $B$, $C$ non alignés. Tracer en couleurs
différentes~: le \omterm{def:g6:lines:objects}{segment} $[AB]$, la demi-droite $[CA)$, la droite
$(BC)$.
\end{exercise}

\begin{exercise}[$\star$]\label{exo:g6:lines:2}
Vrai ou faux~? «~$[AB]$ et $[BA]$ sont le même \omterm{def:g6:lines:objects}{segment}.~»
«~$[AB)$ et $[BA)$ sont la même demi-droite.~» «~$(AB)$ et $(BA)$
sont la même droite.~» Expliquer chaque réponse.
\end{exercise}

\begin{exercise}[$\star$]\label{exo:g6:lines:3}
Tracer une droite $d$ et un point $P$ hors de $d$. Construire avec
l'équerre~: la \omterm{def:g6:lines:perp}{perpendiculaire} à $d$ par $P$, puis la \omterm{def:g6:lines:perp}{parallèle} à $d$
par $P$. Décrire les étapes.
\end{exercise}

\begin{exercise}[$\star$]\label{exo:g6:lines:4}
Les droites $a$ et $b$ sont toutes deux \omterm{def:g6:lines:perp}{perpendiculaires} à une droite
$c$, et une quatrième droite $e$ est \omterm{def:g6:lines:perp}{perpendiculaire} à $a$. Que
peut-on dire de $a$ et $b$~? De $e$ et $c$~? Justifier avec
\cref{prop:g6:lines:facts}.
\end{exercise}

\begin{exercise}[$\star$]\label{exo:g6:lines:5}
Tracer un cercle de centre $O$ de \omterm{def:g3:shapes:shapes}{rayon} $3$~cm. Placer un point $M$
sur le cercle, un point $N$ à l'intérieur, un point $P$ à
l'extérieur. Que peut-on dire des distances $OM$, $ON$, $OP$ par
rapport à $3$~cm~?
\end{exercise}

\begin{exercise}[$\star$]\label{exo:g6:lines:6}
Un cercle a un \omterm{def:g4:geometry:circle}{diamètre} de $9$~cm. Quel est son \omterm{def:g3:shapes:shapes}{rayon}~? Un autre a un
\omterm{def:g3:shapes:shapes}{rayon} de $2.6$~cm~: quel est son \omterm{def:g4:geometry:circle}{diamètre}~?
\end{exercise}

\begin{exercise}[$\star$]\label{exo:g6:lines:7}
Nommer l'angle marqué en trois lettres, puis le classer (aigu, droit,
obtus, plat)~: un angle de $72^\circ$ au \omterm{def:g5:solids:def}{sommet} $R$ entre les
demi-droites vers $S$ et $T$~; un angle de $148^\circ$ au \omterm{def:g5:solids:def}{sommet} $B$
entre les demi-droites vers $A$ et $C$~; un angle de $90^\circ$ au
\omterm{def:g5:solids:def}{sommet} $O$ entre les demi-droites vers $M$ et $N$.
\end{exercise}

\begin{exercise}[$\star$]\label{exo:g6:lines:8}
Estimer, puis mesurer au rapporteur, les trois angles d'un triangle
dessiné soi-même. Additionner les trois mesures~: que trouve-t-on~?
(Garder sa réponse pour \cref{ch:g7:triangles}.)
\end{exercise}

\begin{exercise}[$\star$]\label{exo:g6:lines:9}
Tracer un angle $\widehat{xOy}$ de $65^\circ$ au rapporteur, puis un
angle de $130^\circ$. Comment obtenir le second à partir du premier
sans le rapporteur~?
\end{exercise}

\begin{exercise}[$\star\star$]\label{exo:g6:lines:10}
Deux villages $A$ et $B$ sont dessinés sur une carte. Où sont les
points situés à $3$~cm de $A$ \emph{et} à $2$~cm de $B$~? Dessiner
une figure montrant combien de tels points il peut y avoir ($0$, $1$
ou $2$ selon la distance $AB$).
\end{exercise}

\begin{exercise}[$\star\star$]\label{exo:g6:lines:11}
Une horloge indique 3~heures~: quel est l'angle entre les deux
aiguilles~? Même question à 5~heures, et à 6~heures. (Un tour complet
est $360^\circ$ pour $12$ heures.)
\end{exercise}

\begin{exercise}[$\star\star\star$]\label{exo:g6:lines:12}
Tracer un \omterm{def:g6:lines:objects}{segment} $[AB]$ de $6$~cm. Construire le point $C$ tel que
$AC = BC = 6$~cm, en n'utilisant que le compas, et mesurer l'angle
$\widehat{CAB}$. Quel triangle a-t-on construit, et que conjecture-t-on
sur ses angles~?
\end{exercise}
